% 本文件由 LaTeX 排版 Agent 根据有效 translation.json 自动生成。
% author=Methodius; work_id=0629; title=关于基督十字架与受难讲道词的三篇残篇

\chapter{关于基督十字架与受难讲道词的三篇残篇}

% structure: p0001 source=h1 target=omitted reason=page-title-already-rendered-by-parent

% structure: p0002 source=h2 target=section reason=relative-heading-rank; base=h2

\section{一}

\Emph{美多德}主教对那些说：\Quote{天主子被钉在十字架上并成为人，这对我们有什么益处呢？祂为何忍受十字架这种刑罚，而非其他刑罚？十字架又有何益处？}\par

基督，天主子，因圣父的命令，与可见受造物交往，为的是推翻暴君——即魔鬼——的统治，将我们的灵魂从它们可怕的束缚中解救出来；由于此，我们整个本性，因邪恶的酒而沉醉，充满了动荡与混乱，绝不能返回对善与有益之事的记忆。因此，它也更容易被引向偶像，因为邪恶已完全淹没它，并蔓延所有世代，由于因不顺从而临到我们肉身帐幕的改变；直到主基督，借着祂生活和现身的肉身，削弱了快感冲击的力量，那力量是武装起来攻击我们的阴间势力借以奴役我们心智的，并使人类摆脱所有邪恶。因为为此目的，主耶稣既穿上我们的肉身，成为人，又按天主的上智安排被钉在十字架上；为的是借着魔鬼曾骄傲地、虚假地自称为神的肉身（它们曾用诡计将我们的灵魂掳至死亡），甚至借着这同一肉身，它们能被推翻，并被发现不是神。因为祂借着成为人，阻止了它们的狂妄抬得更高；为的是借着那使理性人类疏远对真天主的敬拜并遭受伤害的身体，甚至因这同一身体以不可言喻的方式接纳了智慧圣言，敌人就会被发现是我们灵魂的毁灭者，而非恩人。因为若基督以祂天主性的可畏和祂不可战胜的权能之伟大，削弱了魔鬼的敌对本性，那本不足为奇。但既然这会给他们带来更大的悲伤和折磨（因为他们宁愿被比自己更强的一位所战胜），因此，祂是通过一个人才获得了族类的安全；为的是使人类，在生命和真理本身以肉身形式进入他们内之后，能够返回圣言的形相和光明，战胜罪恶诱惑的力量；并且使魔鬼，被比自己更弱的一位所征服，因而受到藐视，从而停止他们的过分自信，其地狱般的忿怒被抑制。主要为此，十字架被引入，作为对抗不义的战利品和威慑物，使人从此不再受制于忿怒，在他弥补了因不顺从所遭受的失败，并合法地战胜了阴间权势，且借着天主的恩赐从一切债务中获得释放之后。因此，既然天主首生的圣言这样以正义的甲胄武装了祂寄居其中的人性，祂就（如前所述）藉十字架的形像战胜了奴役我们的权势，并显明了那曾被腐败（如同暴君势力）压迫的人，得以自由，双手不受束缚。因为十字架，你若愿意定义它，就是胜利的确认，天主藉以降至人间的道路，对抗物质神灵的战利品，死亡的击退，上升至真日的根基；为那些急于享受那里之光的人所设的梯子，那使适合教会建造的人从下面被提升、如同四方的石头被嵌合在神圣圣言上的机械。因此，我们的君王觉察到十字架的形像用于驱散一切邪恶，便制作了\OriginalTerm{军旗}{vexillas}，正如在拉丁语中所称呼的那样。因此，海洋顺从这形像，便使自己可为人航行。因为可以说，每个受造物，为了自由，都被这标记所烙印；因为高飞的鸟儿，通过伸展翅膀，形成十字架的形像；而人自己，也以双手伸展，代表同一形像。因此，当主以这形像塑造了他（即从起初祂塑造他的形像），祂将他的身体与天主性结合，为使其从此成为奉献给天主的工具，摆脱一切不和谐与失调。因为人，在他被塑造为敬拜天主、并好似唱过不朽坏的真道之歌、并因此被造为能持有天主性、被装配为生命的竖琴的弦与琴弦之后，人不能（我说）再回到不和谐与腐败。\par

% structure: p0005 source=h2 target=section reason=relative-heading-rank; base=h2

\section{二}

\Emph{同一位美多德致那些以基督十字架为耻的人}\par

有人也以为天主，就是他们用自己的情感尺度去衡量的那一位，也判断恶人和愚人所判断为可赞美或可指责的事，并以人的意见为祂的规则和尺度，却不考虑：由于他们内在的无知，每个受造物都达不到天主的美。因为祂藉着自己的圣言，从万物的普遍本体和本性中，吸引一切归向生命。因为无论祂意愿善，祂自己就是至善，并住在自己内；或者，如果美悦乐祂，既然祂自己就是唯美者，祂便注视自己，丝毫不看重那些令人类赞叹的事。确实，那才应被视为真正最美和最可赞美的，就是天主自己所认为是美的，即使它被其余一切所轻看和鄙视——而不是人所幻想为美的。因此，虽然祂藉此形像愿意将灵魂从败坏的情欲中解救出来，使魔鬼们蒙受显著的羞辱，我们却应当接受它，而不应毁谤它，因为它是赐予我们，为解救我们，并使我们从因不服从而招致的锁链中获得自由的。因为圣言受难，带着肉身被钉在十字架上，为的是将曾被错误欺骗的人带到祂至高而神圣的尊威前，使他恢复到自己所疏远的那个神圣生命。藉此形像，实在，情欲被削弱了；因受难而发生了情欲的情欲，因基督的死亡而发生了死亡的死亡，祂既未被死亡制服，也未被受难的痛苦所胜过。因为受难并未将祂从祂的平静中击倒，死亡也未曾伤害祂，但祂在可受感的（状态）中仍然是不受感的，在有死的（状态）中仍然是不死的，包含空气、这中间状态和天上的一切，并调和有死之于不死的神性。死亡完全被征服了；肉身被钉在十字架上，以引出它的不死性。\par

% structure: p0008 source=h2 target=section reason=relative-heading-rank; base=h2

\section{3}

\Emph{同一位\Person{梅多迪乌斯}：基督天主子如何在短暂而确定的时间内被身体所包围，且原本不受感，却成了受难的牺牲。}\par

因为既然这德能在祂内，那么，能力的本质就在于能在狭小空间内收缩并减少，又能在广阔空间内扩展并增加。但如果祂能与大的扩展且成为相等，却不能与小的收缩并减少，那么能力就不在祂内。因为如果你说这对于能力是可能的，而对于那个是不可能的，你就否定了它是能力；因为它在不能做的事上就是软弱无能的。再者，它也不会在那些经历变化的事上包含任何神性的卓越。因为人和其他动物，对于他们能成就的事，就发挥力量；但对于他们不能实行的事，就软弱而衰亡。因此，天主子被封闭在人性中，就是因为这对祂来说并非不可能。因为祂以大能受苦，仍是不受感的；祂死亡，却将不死的恩赐赐予凡人。因为身体被另一个身体击打或切割时，只被击打到击打者所击打的限度，或切割者所切割的限度。因为根据被击打物的反弹，击打会反射到击打着，因为行动者和承受者两者必须同等承受。如果被切割之物因体积极小，与切割者不相称，就完全不能切割它。因为如果受作用的身体不抵抗剑的击打，反而顺服，那么作用便会无效，如同人们在火与空气的稀薄精细身体上所见的；因为在这种情况下，较坚固物体的冲力被减弱，而徒劳无功。但如果火、空气、石头、铁，或人类为相互毁灭而用来对付彼此的任何东西——如果由于它们具有精微的本性，不能穿透或分裂它们，那么智慧为何不更应该不受伤害且不受感的，丝毫无损于任何事物，即便它与那被钉穿孔的身体结合，因为它比任何其他本性都更纯洁和卓越，如果你只除开那生祂的天主的本性呢？\par

\SourceRecord{Translated by William R. Clark.；Ante-Nicene Fathers；Vol. 6.；Edited by Alexander Roberts, James Donaldson, and A. Cleveland Coxe.；Buffalo, NY: Christian Literature Publishing Co.,；1886.；<http://www.newadvent.org/fathers/0629.htm>.}
